% Source: Wald GR, physical PDF pp. 68--72
%         (printed pp. 55--59).
% Coverage: Section 4.1, The Geometry of Space in Prerelativity Physics;
%           General and Special Covariance; equations (4.1.1)--(4.1.5).
% Figures/tables/footnotes: footnote 1; no figures or tables.
% Status: source-reviewed and compile-clean.
% Uncertainties: none.

\subsection{The Geometry of Space in Prerelativity Physics; General and Special
\mbox{Covariance}}\label{sec:4.1}

In prerelativity physics it is assumed that space has the manifold structure of
\(\mathbb R^3\). It is further assumed that the association of points of space
with elements \((x^1,x^2,x^3)\) of \(\mathbb R^3\) can be achieved in a
natural manner by construction of a \textquotedblleft rigid rectilinear
grid\textquotedblright{} of metersticks. The coordinates of space obtained in
this manner are referred to as \emph{Cartesian coordinates}. Many different
systems of rigid grids (i.e., many Cartesian coordinate systems) are
possible---specifically they can be put in one-to-one correspondence with
elements of the six-parameter group of rotations and translations of
\(\mathbb R^3\). Thus, the Cartesian coordinates \((x^1,x^2,x^3)\) of a
point in space do not, in themselves, have any intrinsic meaning. However, the
distance, \(D\), between two points, \(x\) and \(\bar x\), defined in terms
of Cartesian coordinates by
\begin{equation}
  D^2=(x^1-\bar x^1)^2+(x^2-\bar x^2)^2+(x^3-\bar x^3)^2
  \tag{4.1.1}\label{eq:4.1.1}
\end{equation}
is independent of the choice of Cartesian coordinate system and thus can be
viewed as describing an intrinsic property of space.

This formula for the distance between two points gives rise to a metric of space
\(h_{ab}\) (in the sense of section 2.3) in the following manner. According to
equation~\eqref{eq:4.1.1}, the distance between two
\textquotedblleft nearby\textquotedblright{} points is
\begin{equation}
  (\delta D)^2=(\delta x^1)^2+(\delta x^2)^2+(\delta x^3)^2 .
  \tag{4.1.2}\label{eq:4.1.2}
\end{equation}
This suggests that the metric of space is given by (in the notation of
equation~(2.3.10))
\begin{equation}
  \dd s^2=(\dd x^1)^2+(\dd x^2)^2+(\dd x^3)^2 ,
  \tag{4.1.3}\label{eq:4.1.3}
\end{equation}
or, in the index notation, in the Cartesian coordinate basis, we have
\begin{equation}
  h_{ab}=\sum_{\mu,\nu}h_{\mu\nu}(\dd x^\mu)_a(\dd x^\nu)_b
  \tag{4.1.4}\label{eq:4.1.4}
\end{equation}
with \(h_{\mu\nu}=\diag(1,1,1)\). This definition of \(h_{ab}\) is
independent of choice of Cartesian coordinate system; i.e., the same tensor field
\(h_{ab}\) would result if we had chosen a different Cartesian system.

Let us examine the geometry determined by equation~\eqref{eq:4.1.4}. Since the
components of the metric in the Cartesian coordinate basis are constants, the
ordinary derivative operator of this coordinate system satisfies
\begin{equation}
  \partial_a h_{bc}=0 .
  \tag{4.1.5}\label{eq:4.1.5}
\end{equation}
Consequently, this ordinary derivative is the derivative operator associated with
\(h_{ab}\), and thus \(\Gamma^a{}_{bc}\) vanishes for this coordinate system
(as also can be seen directly from equation~(3.1.30)). Since ordinary derivatives
commute on all tensors, by equation~(3.2.3) the curvature vanishes, i.e., the
metric \(h_{ab}\) is flat. By equation~(3.3.5), the geodesics of space are
precisely those curves which are \textquotedblleft straight
lines\textquotedblright{} when expressed in Cartesian coordinates, i.e., the
curves whose Cartesian coordinates are linearly related to the affine parameter.
Consequently, there is a unique geodesic connecting any given pair of points, and
the length, equation~(3.3.7), of this geodesic is given by
equation~\eqref{eq:4.1.1}. Thus, our metric, equation~\eqref{eq:4.1.4},
reproduces the distance formula which motivated its definition.

Thus, our assumptions about space in prerelativity physics have led to the
statement that space is the manifold \(\mathbb R^3\) which possesses a flat
Riemannian metric. Conversely, given that space is \(\mathbb R^3\) with a flat
Riemannian metric, we can derive all our initial assumptions. We can use the
geodesics of the flat metric to construct a Cartesian coordinate system, using the
fact that initially parallel geodesics remain parallel because the curvature
vanishes. The distance formula, equation~\eqref{eq:4.1.1}, will hold, and hence
one may lay down metersticks over the Cartesian coordinate lines to construct a
\textquotedblleft rigid rectilinear grid\textquotedblright. Thus, everything we
have said thus far in this section is encapsulated in the statement that
\emph{space is the manifold \(\mathbb R^3\) with a flat Riemannian metric
defined on it}.

Let us consider, now, quantities of physical interest in space. All experiments in
physics measure numbers, so all quantities of physical interest must eventually be
reducible to numbers. However, many quantities of interest---such as the magnetic
field or the stress tensor mentioned at the beginning of section 2.3---require the
additional specification of a basis of vectors in order to produce numbers. A very
general class of quantities of interest are maps of vectors and covectors into
numbers. Since any such analytic map can be Taylor-expanded as a sum of
multilinear maps, we see that tensor fields---i.e., multilinear maps of vectors and
covectors into numbers---encompass an extremely wide class of quantities. Indeed,
the generality of the mathematical notion of tensor fields appears to be great
enough so that essentially all\footnote{Spinor fields arise in physics and comprise
a more general class of objects than tensor fields; see chapter 13.} quantities
which one considers in physics can be viewed as tensor fields. The laws of physics
governing these quantities can be expressed as tensor equations, i.e., equalities
between tensor fields defined on space.

An important principle---which goes under the name of \emph{general
covariance}---applies to the form of the laws of physics in the prerelativity
description of space as well as in special relativity and general relativity. The
principle of general covariance in this context states that the metric of space is
the only quantity pertaining to space that can appear in the laws of physics.
Specifically, there are no preferred vector fields or preferred bases of vector
fields pertaining only to the structure of space which appear in any law of
physics. This idea played an important role in motivating the formulation of
special relativity and general relativity, but the principle is rather vague because
the phrase \textquotedblleft pertaining to space\textquotedblright{} does not have
a precise meaning.

Historically, there has been a great deal of discussion concerning the tensor
nature of the laws of physics and the principle of general covariance, so it is
worthwhile to remark here upon other formulations of these ideas. In many
treatments, it is assumed that a coordinate system has been chosen and the
equations of physics have been written out in component form using the coordinate
basis. (Note that, in our viewpoint, this already assumes that the laws of physics
are tensor equations.) Now, if our formulation of general covariance were
violated, say by the existence of a preferred vector field \(v^a\), it would be
possible to adapt a coordinate system to this vector field so that, say,
\((\partial/\partial x^1)^a\) equals \(v^a\). If we wrote out the components
of an equation of physics in such an adapted coordinate system
\emph{without explicitly incorporating \(v^a\) into the equation} but rather
substituting everywhere the components \(v^\mu=(1,0,\ldots,0)\), we would
find, of course, that the form of the equation is not preserved when we make a
coordinate transformation which violates our condition
\((\partial/\partial x^1)^a=v^a\). Thus in these treatments it would be
concluded that in this example the equations are not preserved under general
coordinate transformations. Furthermore, it would be concluded that the equations
are not tensor equations because the components fail to transform according to
equation~(2.3.8) under coordinate transformations. However, in our viewpoint,
the \textquotedblleft nontensorial\textquotedblright{} nature of the equations can
be attributed to a failure to explicitly incorporate the extra geometrical structure
into the equation. When this is done, the equation will have a tensorial character,
but our formulation of the principle of general covariance will be violated by the
appearance of additional quantities pertaining to space.

A good example of an implication of the principle of general covariance which
illustrates this difference in viewpoint is the following statement: A Christoffel
symbol \(\Gamma^c{}_{ab}\) cannot appear (by itself, in undifferentiated form) in
any law of physics. From our viewpoint, \(\Gamma^c{}_{ab}\) is a tensor, but it
is equivalent to specifying an ordinary derivative, \(\partial_a\). However, this
ordinary derivative is an additional geometric quantity pertaining to space which
is not derivable from the metric (unless it coincides with the derivative operator
satisfying equation~(3.1.22), in which case \(\Gamma^c{}_{ab}=0\)), so the
appearance of \(\Gamma^c{}_{ab}\) violates our formulation of general
covariance. From the other viewpoint, \(\Gamma^c{}_{ab}\) cannot appear
because, as already remarked in section 3.1, its coordinate components do not
transform according to equation~(2.3.8) under general coordinate
transformations.

The laws of prerelativity physics also obey the principle of \emph{special
covariance}. The meaning of this principle can be explained as follows. The
metric of space, equation~\eqref{eq:4.1.4}, has a nontrivial group of isometries
(see appendix C), namely the six-parameter group of rotations and translations of
\(\mathbb R^3\) together with the discrete parity symmetry. Consider a family,
\(O\), of observers stationed in space who make measurements on a physical
field. Consider another family, \(O'\), of observers obtained by
\textquotedblleft acting\textquotedblright{} on \(O\) with an isometry. The
principle of special covariance says that any physically possible set of
measurements obtained by \(O\) also is a physically possible set of measurements
for \(O'\). It implies the existence of an action of the isometry group on the
states of the physical fields being measured. The importance of special covariance
lies mainly in the fact that in appropriate contexts, one can use this group action
to derive, or at least motivate, the possible laws governing the physical fields.
This will be illustrated in chapter 13, where---in the context of special relativity
rather than prerelativity physics---special covariance will be used to motivate the
dynamical laws satisfied by fields of mass \(m\) and spin \(s\).

The principle of special covariance is closely related to the principle of general
covariance. Suppose an object of physical interest is described by a tensor field
\(T^{a\cdots}{}_{b\cdots}\). If general covariance holds, then the equations
governing \(T^{a\cdots}{}_{b\cdots}\) should involve only
\(T^{a\cdots}{}_{b\cdots}\), the metric \(h_{ab}\), and the quantities
determined by \(h_{ab}\), such as its derivative operator. Furthermore, all
quantities measurable by \(O\) should be expressible as scalars resulting from
contracting \(T^{a\cdots}{}_{b\cdots}\) and its derivatives with the basis
vector fields \((e_\alpha)^a\) associated with \(O\). Let \(\phi\) be an
isometry. Then if the above assumptions hold, the action of \(\phi\) on
\(T^{a\cdots}{}_{b\cdots}\), defined in appendix C, will produce a tensor field
\(\phi^*T^{a\cdots}{}_{b\cdots}\) which will satisfy the equations for the
physical field and will yield the same measurable quantities for \(O'\) as
\(T^{a\cdots}{}_{b\cdots}\) yields for \(O\). Thus, for physical quantities
describable by tensor fields and satisfying tensor equations, special covariance of
the physical laws under isometries follows, in essence, from general covariance.
(However, note that the formulation of special covariance given above does not
require that the physical quantities be describable as tensor fields, and, indeed,
in chapter 13 we will use special covariance to motivate the introduction of spinor
fields.) On the other hand, the principle of general covariance may be viewed as
being, essentially, a formulation of the idea of special covariance which is
applicable in the absence of isometries.

Special covariance of the laws of physics for tensor fields implies that if we pick
a coordinate system and write out the coordinate component equations
\emph{without explicitly incorporating the metric into the equation} but rather
substituting everywhere its coordinate components \(h_{\mu\nu}\) [e.g.,
\(\diag(1,1,1)\) in Cartesian coordinates], then the form of the equations will
be preserved under the coordinate transformations corresponding to the
isometries, for the simple reason that the numerical values of the components
\(h_{\mu\nu}\) remain unchanged under these transformations. Thus, special
covariance can be viewed as expressing the invariance of the component form of
equations under a \textquotedblleft special\textquotedblright{} group of
coordinate transformations (namely, those corresponding to isometries), while
general covariance can be viewed as expressing a (quite different type of)
invariance of the equations under general coordinate transformations. Historically
these formulations of special and general covariance played a large role in
discussions of relativity theory, and, indeed, the names
\textquotedblleft special\textquotedblright{} and \textquotedblleft
general\textquotedblright{} relativity derive from these formulations of special
and general covariance.
